\section*{ID7883: XI.3 Bessel-Mode Imprinting on the Condensate}
\paragraph{Metadata.}
PiID: 8776182582976 | RTEID: RTE:P30611.4396:T0.1728 | UUID: 8033b3dc-5ad7-587a-a337-22376fdd1472

\paragraph{Canonical Statement.}
\textbackslash{}[ \textbackslash{}phi(\textbackslash{}mathbf\{x\},t) \textbackslash{}approx \textbackslash{}sum\_\{n,m\} A\_\{n,m\}(t) J\_m(k\_\{n,m\} r) e\textasciicircum{}\{i m \textbackslash{}theta\}, \textbackslash{}]

\paragraph{Canonical Equation.}
\[
\phi(\mathbf{x},t) \approx \sum_{n,m} A_{n,m}(t) J_m(k_{n,m} r) e^{i m \theta},
\]

\paragraph{Dictionary.}
XI.3 Bessel-Mode Imprinting on the Condensate — φ(x,t) ≈ ∑\_n,m A\_n,m(t) J\_m(k\_n,m r) e\textasciicircum{}i m θ

\paragraph{Encyclopedia.}
XI.3 Bessel-Mode Imprinting on the Condensate is a series / summation, written φ(x,t) ≈ ∑\_n,m A\_n,m(t) J\_m(k\_n,m r) e\textasciicircum{}i m θ. It is expressed as a sum or series over modes or indices, superposing discrete contributions. It is built from φ, θ, x, n, m, A\_n,m. Within the Trawin Zero Point registry it serves as one element of the framework's mathematical narrative.
